%===========================================================
% Capítulo: Conclusiones y trabajo futuro — Síntesis física, estadística y clínica
%===========================================================
\chapter{Conclusiones}
\label{chap:conclusiones}

El trabajo demostró que la clasificación directa de hologramas en línea de eritrocitos, sin reconstrucción numérica de amplitud o fase, es una ruta técnicamente sólida y clínicamente prometedora para el tamizaje de anemia falciforme en plataformas de microscopía holográfica digital sin lente (LDHM) de bajo costo. La piedra angular del enfoque es la proyección físico–estadística de la señal interferométrica hacia un espacio de rasgos parco pero expresivo que conserva la información de escala, periodicidad y direccionalidad codificada por la difracción en régimen escalar de Fresnel. Al anclar los descriptores en magnitudes invariantes y fisi\-camente interpretables —el perfil radial del espectro en fracción de Nyquist, la localización de picos asociados a la periodicidad de franjas, la razón de energías en bandas que cuantifica el enriquecimiento de altas frecuencias y la anisotropía angular que condensa la elipticidad del patrón—, se evita la dependencia de calibraciones finas de distancia objeto–sensor o del filtrado de imagen gemela, preservando además la robustez frente a variaciones suaves de fotometría y ruido. Esta síntesis otorga sentido físico a la discriminación estadística: los eritrocitos falciformes, al perder circularidad y ganar rigidez y elongación, inducen franjas más juntas y direccionalidad marcada en el holograma crudo, por lo que desplazan hacia arriba los picos radiales y elevan la energía espectral relativa en bandas medias–altas, a la vez que aumentan la entropía y el contraste de texturas locales, todo ello coherente con el modelo de fase delgada y con la estructura de interferencia esperable.

Desde la perspectiva metodológica, un aporte central es la evaluación honesta del rendimiento mediante validación cruzada anidada estratificada por sujeto y el uso de \emph{bootstrap} a nivel de sujeto para cuantificar incertidumbre. La separación estricta entre selección de hiperparámetros y estimación, junto con la decisión reportada en la granularidad clínicamente relevante, reduce el sesgo optimista habitual en estudios de pocos datos y fortalece la reproducibilidad. La calibración probabilística sobre predicciones \emph{out-of-fold} permitió traducir los puntajes a probabilidades compatibles con curvas de confiabilidad, Brier score y análisis de utilidad clínica basados en umbrales de costo; en escenarios de cribado, donde el costo del falso negativo domina, el sistema sostuvo sensibilidades elevadas con especificidades razonables, generando beneficios clínicos netos frente a estrategias triviales en rangos operativos realistas. Igualmente relevante es la parquedad alcanzada: al reducir el vector de rasgos a un núcleo compacto sin sacrificar AUC ni AUPRC más allá de márgenes mínimos, se mejora la estabilidad inter–pliegue, se facilita la auditoría y se clarifica la narrativa explicativa post–hoc, en la que los valores de Shapley ubican consistentemente al primer pico radial, la razón espectral en alta–baja y los descriptores de textura como motores de la decisión.

En términos de ingeniería, la cadena completa —normalización por percentiles, CLAHE controlado, estandarización $z$-score, muestreo determinista de parches, extracción de rasgos en coordenadas físicas y agregación robusta parche→sujeto— se implementó con encapsulamiento suficiente para prevenir fugas y con costos de cómputo compatibles con CPU de propósito general. Este balance entre fundamento físico, corrección estadística y eficiencia computacional habilita el tránsito hacia prototipos de punto de atención (\emph{POCT}) sobre dispositivos económicos, con latencias por sujeto del orden de décimas de segundo y requerimientos de memoria modestos, sin demandar reconstrucciones ni redes profundas de gran tamaño. La coherencia de los hallazgos a través de estratos de SNR y pequeñas perturbaciones fotométricas sugiere que el modelo capta estructura interferométrica genuina y no artefactos de preprocesamiento, reforzando su viabilidad para integrarse en flujos clínicos de cribado que priorizan rapidez, costo y portabilidad.

No obstante, la validez externa permanece como límite explícito de inferencia. La ausencia de una cohorte verdaderamente \emph{out-of-domain} impide reclamar generalización fuerte entre laboratorios con instrumentación y preparación de muestra distintas; aunque expresar las frecuencias en fracción de Nyquist, utilizar bandas log–espaciadas y apelar a descriptores multiescala amortigua desplazamientos suaves, la confirmación independiente en nuevos sitios es necesaria para consolidar la utilidad del sistema. De igual forma, si bien el manejo de la imagen gemela se resolvió pragmáticamente al operar sobre la intensidad interferométrica y diseñar rasgos robustos a superposiciones, podrían existir condiciones de adquisición que exacerben aliasing angular o patrones parásitos con firmas espectrales similares a las celulares, por lo que políticas de control de calidad y repetición de captura son recomendables cuando la SNR espectral se ubica consistentemente en cuartiles bajos. Por último, el tamaño muestral y la distribución de prevalencia condicionan el ancho de los intervalos de confianza: aun con AUC altas, IC amplios requieren comunicar la incertidumbre de manera transparente y adoptar umbrales operativos conservadores en despliegues preliminares.

En conjunto, la evidencia presentada sustenta la tesis de que el análisis directo de hologramas crudos con una proyección físico–estadística parca puede constituir el núcleo de un tamizaje portátil para anemia falciforme, con fundamentación óptica clara, protocolo estadístico correcto y desempeño compatible con decisiones clínicas en entornos de recursos limitados. El trabajo ofrece, además, una plantilla metodológica reutilizable para otras tareas de citometría holográfica en línea, donde la reconstrucción puede ser sustituida por descriptores informativamente suficientes y verificables contra la física subyacente.

\section*{Trabajo futuro y proyección}
Las líneas de continuación más inmediatas se orientan a robustecer validez externa, ampliar interpretabilidad y consolidar el flujo operativo en \emph{POCT}. En primer lugar, se propone evaluar el sistema en una cohorte externa con variaciones deliberadas de tamaño de píxel, distancia objeto–sensor y preparación de muestra, adoptando un protocolo de \emph{transfer learning} mínimo centrado en recalibración de umbral y, si fuese necesario, en un ajuste isotónico con un pequeño lote de validación local; este esquema preserva el subespacio de rasgos físicos, limita grados de libertad y privilegia la reproducibilidad. En paralelo, la incorporación de una medida de coherencia anular que compare la energía por banda con su varianza radial normalizada permitiría penalizar patrones de alta frecuencia dispersa generados por suciedad o rugosidad de portaobjetos, separándolos de anillos celulares organizados; esta métrica, de inspiración óptica, podría integrarse sin alterar la parquedad global. Otra vía compatible con el presupuesto computacional es un clasificador híbrido que mantenga los rasgos físicos como eje y añada una proyección aprendida de muy baja dimensión con regularización fuerte (norma nuclear o \emph{dropout} pesado), de manera que cualquier ganancia atribuible a interacciones de orden superior se demuestre con validación anidada y estimación honesta de incertidumbre.

Más allá del rendimiento, la agenda incluye una capa de monitorización de calidad que exponga indicadores de SNR espectral, saturación y coherencia angular por sujeto, junto con recomendaciones automáticas de reacquisición cuando la señal cae por debajo de umbrales preestablecidos. Esta \emph{telemetría} integrada mejora la seguridad del sistema y reduce la probabilidad de decisiones en condiciones no diagnósticas. Desde el punto de vista de regulación y ética, conviene preparar documentación de ciclo de vida del modelo, con trazabilidad de datos, versiones, semillas y artefactos, además de un plan de vigilancia post–despliegue que monitoree \emph{drift} de distribución y rendimiento mediante pruebas de \emph{canary} y auditorías periódicas. Finalmente, una extensión natural del marco es la detección multi–clase de morfologías eritrocitarias patológicas (esferocitosis, eliptocitosis, drepanocitosis) a partir de la misma señal interferométrica; la desagregación de firmas angulares por modos periódicos y la separación más fina de picos radiales —eventualmente con \emph{super–resolution} espectral— abren la puerta a un clasificador jerárquico interpretable que conserve la filosofía de no reconstrucción y potencie el valor clínico del LDHM en hematología.

\section*{Implicaciones para salud pública y adopción en campo}
El eje práctico de este trabajo es su compatibilidad con dispositivos económicos y flujos de trabajo rápidos. En entornos donde el acceso a HPLC o electroforesis es limitado, un tamizaje sensible de bajo costo puede reducir tiempos de derivación y pérdidas de seguimiento, incrementando la proporción de lactantes y niños que alcanzan diagnóstico confirmatorio y manejo temprano. La portabilidad del pipeline, su transparencia explicativa y su capacidad de reporte con probabilidades calibradas y métricas con intervalos de confianza facilitan la apropiación por equipos clínicos y de laboratorio, a la vez que sientan bases para su integración con sistemas de registro y programas de cribado neonatal. El paso de viabilidad técnica a impacto poblacional exige, sin embargo, estudios de efectividad en condiciones reales, con medición de desenlaces operativos (tiempos, costos, cobertura) y clínicos (oportunidad de diagnóstico, inicio de profilaxis, hospitalizaciones evitadas), de modo que las ventajas demostradas aquí en discriminación y eficiencia computacional se traduzcan en beneficios tangibles.

\section*{Cierre}
El valor del enfoque no descansa sólo en que evita la reconstrucción, sino en que lo hace sin renunciar a la ciencia de la formación de imagen: las magnitudes que gobiernan la decisión son observables y previsibles a partir del modelo de difracción, lo que eleva el estándar de explicabilidad en visión biomédica. En un campo a menudo dominado por cajas negras, esta conjunción de óptica, estadística y diseño de sistemas representa un camino fértil para tecnologías de diagnóstico accesibles, auditables y efectivas.

%===========================================================
% (En el siguiente bloque se presentarán: Resumen, Abstract y Anexos sugeridos con figuras/tablas.)
%===========================================================
